\exercisesection{\S~8}

\begin{enumerate}
\def\labelenumi{\arabic{enumi})}
\item
  Let \(f\) be an invariant polynomial function on \(\mathfrak{g}\) . Show that \(f\) is invariant under \(\operatorname{Aut}(\mathfrak{g})\) if and only if \(f|\mathfrak{h}\) is invariant under \(\operatorname{Aut}(\mathrm{R})\) . Deduce that, if the Dynkin graph of \(\mathrm{R}\) has a non-trivial automorphism, there exists an invariant polynomial function on \(\mathfrak{g}\) that is not invariant under \(\operatorname{Aut}(\mathfrak{g})\) .
\item
  Take \(\mathfrak{g} = \mathfrak{sl}(3,k)\) . Show that \(x\mapsto \det (x)\) is an invariant polynomial function on \(\mathfrak{g}\) that is not invariant under \(\operatorname {Aut}(\mathfrak{g})\) (use the automorphism \(x\mapsto -^t x\) ).
\item
  Let \(\mathfrak{a}\) be a semi-simple Lie algebra, and \(s\in \mathrm{Aut}(\mathfrak{a})\) . Show the equivalence of:
\end{enumerate}

\begin{enumerate}
\def\labelenumi{(\roman{enumi})}
\item
  \(s \in \text{Aut}_{0}(\mathfrak{a})\) .
\item
  \(s\) operates trivially on the centre of \(\mathrm{U}(\mathfrak{a})\) .
\item
  For all \(x \in \mathfrak{a}\) , there exists \(t \in \mathrm{Aut}_0(\mathfrak{a})\) such that \(tx = sx\) .
\end{enumerate}

(Use Prop. 6 to show that (iii) \(\Longrightarrow\) (i).)

\begin{enumerate}
\def\labelenumi{\arabic{enumi})}
\setcounter{enumi}{3}
\item
  Show that, in Cor. 2 of Prop. 2, and in Th. 1 (ii), we can restrict ourselves to representations \(\rho\) whose weights are radical weights (remark that Prop. 1 remains valid when \(k[\mathrm{P}]^{\mathrm{W}}\) is replaced by \(k[\mathrm{Q}]^{\mathrm{W}}\) , where \(\mathrm{Q}\) is the group of radical weights).
\item
  We retain the notations of §6 and §7. Let \(\lambda \in \mathfrak{h}^*\) . If, for any \(w \in \mathrm{W}\) , \(w \neq 1\) , we have \((\lambda + \rho) - w(\lambda + \rho) \notin \mathrm{Q}_+\) , then \(\mathrm{Z}(\lambda)\) is simple.
\end{enumerate}

(Use Cor. 1 (ii) of Th. 2.)

\begin{enumerate}
\def\labelenumi{\arabic{enumi})}
\setcounter{enumi}{5}
\item
  Let \(\mathfrak{a}\) be a Lie algebra, \(f\) a polynomial function on \(\mathfrak{a}\) , and \(x, y\) two elements of \(\mathfrak{a}\) . Put \(f_y = \theta^*(y)f\) (cf.~no. 3), and denote by \(\mathrm{D}_x f\) the tangent linear map of \(f\) at \(x\) (Chap. VII, App. I, no. 2). Show that \(f_y(x) = (\mathrm{D}_xf)([x,y])\) . Deduce that \(\mathrm{D}_xf\) vanishes on \(\operatorname{Im}(\operatorname{ad}x)\) when \(f\) is invariant.
\item
  *Let \(d_1, \ldots, d_l\) be the characteristic degrees of the algebra I(g), cf.~Chap. V, §5, no. 1. For any integer \(n \geq 0\) , denote by \(r_n\) the number of elements of degree \(n\) in a homogeneous basis of S(g) over I(g) (cf.~no. 3, Remark 2), and put \(r(\mathrm{T}) = \sum_{n=0}^{\infty} r_n \mathrm{T}^n\) . Show that
\end{enumerate}

\[
r (\mathrm{T}) = (1 - \mathrm{T}) ^ {- \mathrm{N}} \prod_ {i = 1} ^ {l} (1 - \mathrm{T} ^ {d _ {i}}), \quad \text { where } \mathrm{N} = \dim (\mathfrak {g}).
\]

\begin{enumerate}
\def\labelenumi{\arabic{enumi})}
\setcounter{enumi}{7}
\tightlist
\item
  Put \(l = \operatorname{rk}(\mathfrak{g}), N = \dim(\mathfrak{g})\) . If \(x \in \mathfrak{g}\) , define \(a_i(x)\) , \(0 \leq i \leq N\) , by the formula
\end{enumerate}

\[
\det (\mathrm{T} + \operatorname{ad} x) = \sum_ {i = 0} ^ {\mathrm{N}} \mathrm{T} ^ {\mathrm{N} - i} a _ {i} (x).
\]

The function \(a_{i}\) thus defined is homogeneous polynomial of degree i, and invariant under \(\operatorname{Aut}(\mathfrak{g})\) . If \(x \in h\) and \(i \leq N - l\) , \(a_{i}(x)\) is the ith elementary symmetric function of the \(\alpha(x)\) , \(\alpha \in R\) ; in particular, \(a_{N-l}(x) = \prod_{\alpha \in \mathbb{R}} \alpha(x)\) . Construct an example in which the \(a_{i}\) do not generate the algebra of polynomial functions on g invariant under \(\operatorname{Aut}(\mathfrak{g})\) .

\begin{enumerate}
\def\labelenumi{\arabic{enumi})}
\setcounter{enumi}{8}
\tightlist
\item
  The notations are those of no. 5. Let \(\lambda \in \mathfrak{h}^*\) , \(z \in \mathbf{Z}\) , \(z'\) the image of \(z\) under the principal anti-automorphism of \(\mathrm{U}(\mathfrak{g})\) , and \(w_0\) the element of W that transforms B into \(-\mathrm{B}\) . Show that \(\chi_{\lambda}(z) = \chi_{-w_0\lambda}(z')\) .
\end{enumerate}

(It suffices to prove this for \(\lambda \in \mathrm{P}_{++}\) . Consider the operation of \(z\) on \(\mathrm{E}(\lambda)\) and \(\mathrm{E}(\lambda)^*\) ; use Prop. 11 of §7.)

\begin{enumerate}
\def\labelenumi{\arabic{enumi})}
\setcounter{enumi}{9}
\tightlist
\item
  (In this exercise, and in the following three, we retain the notations of §6.) Let \(\lambda \in \mathfrak{h}^*\) .
\end{enumerate}

\begin{enumerate}
\def\labelenumi{\alph{enumi})}
\tightlist
\item
  Let \(\mathrm{N}, \mathrm{N}'\) be \(\mathfrak{g}\) -submodules of \(\mathrm{Z}(\lambda)\) such that \(\mathrm{N}' \subset \mathrm{N}\) and \(\mathrm{N} / \mathrm{N}'\) is simple. Show that there exists \(\mu \in \lambda - \mathrm{Q}_+\) such that \(\mathrm{N} / \mathrm{N}'\) is isomorphic to \(\mathrm{E}(\mu)\) (apply Th. 1 of §6), and such that \(\mu + \rho \in \mathrm{W}.(\lambda + \rho)\) (apply Cor. 1 of Th. 2).
\item
  Show that \(\mathrm{Z}(\lambda)\) admits a Jordan-Hölder sequence. (Apply a) and the fact that the weights of \(\mathrm{Z}(\lambda)\) are of finite multiplicity.)
\end{enumerate}

\begin{enumerate}
\def\labelenumi{\arabic{enumi})}
\setcounter{enumi}{10}
\tightlist
\item
  Let \(\lambda \in \mathfrak{h}^*\) and let \(V\) be a non-zero \(\mathfrak{g}\) -submodule of \(Z(\lambda)\) . Show that there exists \(\mu \in \mathfrak{h}^*\) such that \(V\) contains a simple \(\mathfrak{g}\) -submodule isomorphic to \(Z(\mu)\) .
\end{enumerate}

(Let A be the set of \(\nu \in \mathfrak{h}^*\) such that V contains a \(\mathfrak{g}\) -submodule isomorphic to \(Z(\nu)\) . By using Prop. 6 of no. 6, show first that \(A \neq \varnothing\) . Then show that A is finite, and consider an element \(\mu\) of A such that \((\mu - Q_{+}) \cap A = \{\mu\}\) .)

¶ 12) a) Let \(\lambda, \mu \in \mathfrak{h}^*\) . Every non-zero \(\mathfrak{g}\) -homomorphism from \(Z(\mu)\) to \(Z(\lambda)\) is injective. (Use Prop. 6 of no. 6.)

\begin{enumerate}
\def\labelenumi{\alph{enumi})}
\setcounter{enumi}{1}
\item
  Let \(r \in \mathbf{N}\) , A a finite subset of \(\mathbf{N}^r\) , \(m = \mathrm{Card}(\mathrm{A})\) . For all \(\xi \in \mathbf{N}^r\) , let \(s(\xi)\) be the sum of the coordinates of \(\xi\) , and \(\mathfrak{P}_{\mathrm{A}}(\xi)\) the number of families \((n_{\alpha})_{\alpha \in \mathrm{A}}\) of integers \(\geq 0\) such that \(\xi = \sum_{\alpha \in \mathrm{A}} n_{\alpha}\alpha\) . Then \(\mathfrak{P}_{\mathrm{A}}(\xi) \leq (s(\xi) + 1)^m\) for all \(\xi \in \mathbf{N}^r\) . (Argue by induction on \(m\) .)
\item
  Let \(\lambda, \mu \in \mathfrak{h}^*\) . Show that \(\dim \operatorname{Hom}_{\mathfrak{g}}(\mathrm{Z}(\mu), \mathrm{Z}(\lambda)) \leq 1\) . (Let \(\varphi_1\) and \(\varphi_2\) be non-zero \(\mathfrak{g}\) -homomorphisms from \(\mathrm{Z}(\mu)\) to \(\mathrm{Z}(\lambda)\) . If \(\operatorname{Im}(\varphi_1) = \operatorname{Im}(\varphi_2)\) , \(\varphi_1\) and \(\varphi_2\) are linearly dependent by Prop. 1 (iii) of §6. Assume that \(\operatorname{Im}(\varphi_1) \neq \operatorname{Im}(\varphi_2)\) . If \(\mathrm{Z}(\mu)\) is simple, the sum \(\operatorname{Im}(\varphi_1) + \operatorname{Im}(\varphi_2)\) is direct; deduce that \(\mathfrak{P}(\xi + \lambda - \mu) \geq 2\mathfrak{P}(\xi)\) for all \(\xi \in \mathfrak{h}^*\) , contradicting \(b\) ). In the general case, use Exerc. 11.)
\end{enumerate}

When \(\dim\operatorname{Hom}_{\mathfrak{g}}(\mathrm{Z}(\mu),\mathrm{Z}(\lambda))=1\) , write \(\mathrm{Z}(\mu)\subset\mathrm{Z}(\lambda)\) by abuse of notation. d) Let \(\nu\in\mathfrak{h}^{*}\) . The set of \(\lambda\in\mathfrak{h}^{*}\) such that \(\mathrm{Z}(\lambda-\nu)\subset\mathrm{Z}(\lambda)\) is closed in \(\mathfrak{h}^{*}\) in the Zariski topology.

¶ 13) a) Let \(\mathfrak{a}\) be a nilpotent Lie algebra, \(x \in \mathfrak{a}, n \in \mathbf{N}, p \in \mathbf{N}\) . There exists \(l \in \mathbf{N}\) such that \(x^{l}y_{1} \ldots y_{n} \in \mathrm{U}(\mathfrak{a})x^{p}\) for all \(y_{1}, \ldots, y_{n} \in \mathfrak{a}\) .

\begin{enumerate}
\def\labelenumi{\alph{enumi})}
\setcounter{enumi}{1}
\tightlist
\item
  Let \(\lambda, \mu \in \mathfrak{h}^*\) , and \(\alpha \in \mathrm{B}\) be such that
\end{enumerate}

\[
\mathrm{Z} (s _ {\alpha} \mu - \rho) \subset \mathrm{Z} (\mu - \rho) \subset \mathrm{Z} (\lambda - \rho).
\]

Assume that \(\lambda \in \mathrm{P}\) . Let \(p = \lambda (H_{\alpha})\in \mathbf{Z}\) . Show that:

\(b_{1})\) If \(p \leq 0\) , then \(\mathrm{Z}(\lambda - \rho) \subset \mathrm{Z}(s_{\alpha}\lambda - \rho)\) .

\[
b _ {2}) \text {   If   } p > 0, \text {   then   } \mathrm{Z} (s _ {\alpha} \mu - \rho) \subset \mathrm{Z} (s _ {\alpha} \lambda - \rho) \subset \mathrm{Z} (\lambda - \rho).
\]

(Use \(a\) ), and §6, Cor. 1 of Prop. 6.)

\begin{enumerate}
\def\labelenumi{\alph{enumi})}
\setcounter{enumi}{2}
\tightlist
\item
  Let \(\lambda \in \mathfrak{h}^*\) , \(\alpha \in \mathrm{R}\) , and \(m = \lambda(H_{\alpha})\) . Assume that \(m \in \mathbf{N}\) . Show that
\end{enumerate}

\[
\mathrm{Z} (s _ {\alpha} \lambda - \rho) \subset \mathrm{Z} (\lambda - \rho).
\]

(Prove this first for \(\lambda \in \mathrm{P}\) by using \(b\) ), and then in the general case by using Exerc. 12 \(d\) .) \(^{17}\)

\begin{enumerate}
\def\labelenumi{\arabic{enumi})}
\setcounter{enumi}{13}
\item
  *Let \(\mathfrak{a}\) be a semi-simple Lie algebra, and \(Z(\mathfrak{a})\) the centre of \(U(\mathfrak{a})\) . Show that \(U(\mathfrak{a})\) is a free \(Z(\mathfrak{a})\) -module. (Remark that \(grU(\mathfrak{a})\) is isomorphic to \(S(\mathfrak{a})\) and \(S(\mathfrak{a}^*)\) , and use Remark 2 of no. 3.)*
\item
  Let \(x\) be a diagonalizable element of \(\mathfrak{g}\) (§3, Exerc. 10), and \(y\) a semi-simple element of \(\mathfrak{g}\) such that \(f(x) = f(y)\) for every invariant polynomial function \(f\) on \(\mathfrak{g}\) . Show that there exists \(s \in \mathrm{Aut}_e(\mathfrak{g})\) such that \(sy = x\) . (Remark that ad \(x\) and ad \(y\) have the same characteristic polynomial, cf.~Exerc. 8, hence the fact that y is diagonalizable. Next, reduce to the case in which x and y are contained in h by using Exerc. 10 of §3, and use Th. 1 (i) and Lemma 6 to prove that x and y are conjugate under W.)
\item
  Let \(\mathfrak{a}\) be a semi-simple Lie algebra, \(x\) an element of \(\mathfrak{a}\) , and \(x_{s}\) the semi-simple component of \(x\) . Show that, if \(f\) is an invariant polynomial function on \(\mathfrak{a}\) , then \(f(x) = f(x_{s})\) . (Reduce to the case in which \(f\) is of the form \(x \mapsto \operatorname{Tr} \rho(x)^n\) .)
\item
  Assume that \(k\) is algebraically closed, and put \(\mathrm{G} = \operatorname{Aut}_e(\mathfrak{g})\) . Let \(x, y \in \mathfrak{g}\) . Show the equivalence of:
\end{enumerate}

\begin{enumerate}
\def\labelenumi{(\roman{enumi})}
\item
  The semi-simple components of \(x\) and \(y\) are G-conjugate.
\item
  For every invariant polynomial function \(f\) on \(\mathfrak{g}\) , we have \(f(x) = f(y)\) . (Use Exerc. 15 and 16.)
\end{enumerate}

\begin{enumerate}
\def\labelenumi{\arabic{enumi})}
\setcounter{enumi}{17}
\tightlist
\item
  Let \(\mathfrak{a}\) be a semi-simple Lie algebra, \(l = \mathrm{rk}(\mathfrak{a})\) , I the algebra of invariant polynomial functions on \(\mathfrak{a}\) , and \(\mathrm{P}_1, \ldots, \mathrm{P}_l\) homogeneous elements of I generating I as an algebra. The \(\mathrm{P}_i\) define a polynomial map \(\mathrm{P} : \mathfrak{a} \to k^l\) . If \(x \in \mathfrak{a}\) , denote by \(\mathrm{D}_x\mathrm{P} : \mathfrak{a} \to k^l\) the tangent linear map of P at \(x\) (Chap. VII, App. I, no. 2).
\end{enumerate}

\begin{enumerate}
\def\labelenumi{\alph{enumi})}
\tightlist
\item
  Let \(\mathfrak{h}\) be a Cartan subalgebra of \(\mathfrak{a}\) , and let \(x \in \mathfrak{h}\) . Prove the equivalence of:
\end{enumerate}

\begin{enumerate}
\def\labelenumi{(\roman{enumi})}
\item
  \(D_{x}P|\mathfrak{h}\) is an isomorphism from \(\mathfrak{h}\) to \(k^{l}\) ;
\item
  \(x\) is regular.
\end{enumerate}

(Reduce to the split case. Choose a basis of \(\mathfrak{h}\) , and denote by \(d(x)\) the determinant of the matrix of \(\mathrm{D}_x\mathrm{P}|\mathfrak{h}\) relative to this basis. Show, by means of Prop. 5 of Chap. V, §5, no. 4, that there exists \(c \in k^*\) such that \(d(x)^2 = c \prod_{\alpha \in \mathrm{R}} \alpha(x)\) ,

where \(\alpha\) belongs to the set \(R\) of roots of \((\mathfrak{a},\mathfrak{h})\) .)

If these conditions are satisfied, \(\mathfrak{a} = \mathfrak{h} \oplus \operatorname{Im} \operatorname{ad}(x)\) , and \(\operatorname{Ker} D_x P = \operatorname{Im} \operatorname{ad}(x)\) (use Exerc. 6 to show that \(D_x P\) vanishes on \(\operatorname{Im} \operatorname{ad}(x)\) ).

\begin{enumerate}
\def\labelenumi{\alph{enumi})}
\setcounter{enumi}{1}
\tightlist
\item
  Show that the set of \(x \in a\) such that \(D_{x}P\) is of rank l is a dense open subset of a in the Zariski topology.
\end{enumerate}

